
\documentclass{article}
\usepackage{colortbl}
\usepackage{makecell}
\usepackage{multirow}
\usepackage{supertabular}

\begin{document}

\newcounter{utterance}

\centering \large Interaction Transcript for game `dond', experiment `coop', episode 1 with claude{-}opus{-}4{-}8{-}azure.
\vspace{24pt}

{ \footnotesize  \setcounter{utterance}{1}
\setlength{\tabcolsep}{0pt}
\begin{supertabular}{c@{$\;$}|p{.15\linewidth}@{}p{.15\linewidth}p{.15\linewidth}p{.15\linewidth}p{.15\linewidth}p{.15\linewidth}}
    \# & $\;$A & \multicolumn{4}{c}{Game Master} & $\;\:$B\\
    \hline

    \theutterance \stepcounter{utterance}  
    & & \multicolumn{4}{p{0.6\linewidth}}{
        \cellcolor[rgb]{0.9,0.9,0.9}{
            \makecell[{{p{\linewidth}}}]{
                \texttt{\tiny{[P1$\langle$GM]}}
                \texttt{Hrajete kooperativní vyjednávací hru, ve které se musíte dohodnout na tom, jak rozdělit sadu předmětů mezi sebou a druhým hráčem.} \\
\\ 
\texttt{Pravidla:} \\
\texttt{(a) Vy a druhý hráč obdržíte sadu předmětů. Každý z vás také obdrží tajnou funkcí hodnoty, která představuje, jak moc si ceníte jednotlivých typů objektů.} \\
\texttt{(b) Vyměňujete si zprávy s druhým hráčem, abyste se dohodli na tom, kdo dostane které předměty. Můžete poslat maximálně 5 zpráv každý nebo hru předčasně ukončit tím, že kdykoli podáte tajný návrh.} \\
\texttt{(c) Od každého z vás se vyžaduje, aby předložil tajný návrh označující položky, které chcete, formátované v hranatých závorkách takto: "{[}Návrh: <číslo> <název objektu>, <číslo> <název objektu>, <...>{]}"} \\
\texttt{(d) Pokud jsou vaše návrhy komplementární, tj. je dostatek předmětů na splnění obou návrhů, každý hráč obdrží skóre na základě součtu hodnot předmětů, které obdržel. Jinak oba dostanete nula bodů.} \\
\texttt{(e) Vaším cílem je maximalizovat součet vašeho skóre a skóre druhého hráče.} \\
\\ 
\texttt{Pojďme začít.} \\
\\ 
\texttt{Sada dostupných položek je:} \\
\\ 
\texttt{2 pera, 2 citrony, 2 hrnce} \\
\\ 
\texttt{Vaše tajná funkce hodnoty je:} \\
\\ 
\texttt{pero: 2, citron: 0, hrnec: 3} \\
\\ 
\texttt{DŮLEŽITÉ: Vaše zprávy, pokud se nejedná o tajný návrh, jsou předány přímo druhému hráči, takže do nich nezahrnujte žádnou odpověď na pravidla ani text oznamující vaši zprávu. Chcete{-}li podat tajný návrh, použijte uvedený formát. Nepoužívejte hranaté závorky při komunikaci s druhým hráčem, jinak bude interpretován jako váš tajný návrh.} \\
            }
        }
    }
    & \\ \\

    \theutterance \stepcounter{utterance}  
    & \multicolumn{4}{p{0.6\linewidth}}{
        \cellcolor[rgb]{0.9,0.9,0.9}{
            \makecell[{{p{\linewidth}}}]{
                \texttt{\tiny{[P1$\rangle$GM]}}
                \texttt{Ahoj! Pojďme se domluvit tak, aby to bylo výhodné pro nás oba. Pro mě mají největší cenu hrnce, pera mají střední hodnotu a citrony mi nic nedávají. Co je nejcennější pro tebe? Rád ti přenechám věci, které si ceníš nejvíc, ať dohromady získáme co nejvíc bodů.} \\
            }
        }
    }
    & & \\ \\

    \theutterance \stepcounter{utterance}  
    & & & \multicolumn{2}{p{0.3\linewidth}}{
        \cellcolor[rgb]{0.9,0.9,0.9}{
            \makecell[{{p{\linewidth}}}]{
                \texttt{\tiny{[GM$|$GM]}}
                \texttt{continue} \\
            }
        }
    }
    & & \\ \\

    \theutterance \stepcounter{utterance}  
    & & & \multicolumn{2}{p{0.3\linewidth}}{
        \cellcolor[rgb]{0.9,0.9,0.9}{
            \makecell[{{p{\linewidth}}}]{
                \texttt{\tiny{[GM$|$GM]}}
                \texttt{Ahoj! Pojďme se domluvit tak, aby to bylo výhodné pro nás oba. Pro mě mají největší cenu hrnce, pera mají střední hodnotu a citrony mi nic nedávají. Co je nejcennější pro tebe? Rád ti přenechám věci, které si ceníš nejvíc, ať dohromady získáme co nejvíc bodů.} \\
            }
        }
    }
    & & \\ \\

    \theutterance \stepcounter{utterance}  
    & & \multicolumn{4}{p{0.6\linewidth}}{
        \cellcolor[rgb]{0.9,0.9,0.9}{
            \makecell[{{p{\linewidth}}}]{
                \texttt{\tiny{[GM$\rangle$P2]}}
                \texttt{Hrajete kooperativní vyjednávací hru, ve které se musíte dohodnout na tom, jak rozdělit sadu předmětů mezi sebou a druhým hráčem.} \\
\\ 
\texttt{Pravidla:} \\
\texttt{(a) Vy a druhý hráč obdržíte sadu předmětů. Každý z vás také obdrží tajnou funkcí hodnoty, která představuje, jak moc si ceníte jednotlivých typů objektů.} \\
\texttt{(b) Vyměňujete si zprávy s druhým hráčem, abyste se dohodli na tom, kdo dostane které předměty. Můžete poslat maximálně 5 zpráv každý nebo hru předčasně ukončit tím, že kdykoli podáte tajný návrh.} \\
\texttt{(c) Od každého z vás se vyžaduje, aby předložil tajný návrh označující položky, které chcete, formátované v hranatých závorkách takto: "{[}Návrh: <číslo> <název objektu>, <číslo> <název objektu>, <...>{]}"} \\
\texttt{(d) Pokud jsou vaše návrhy komplementární, tj. je dostatek předmětů na splnění obou návrhů, každý hráč obdrží skóre na základě součtu hodnot předmětů, které obdržel. Jinak oba dostanete nula bodů.} \\
\texttt{(e) Vaším cílem je maximalizovat součet vašeho skóre a skóre druhého hráče.} \\
\\ 
\texttt{Pojďme začít.} \\
\\ 
\texttt{Sada dostupných položek je:} \\
\\ 
\texttt{2 pera, 2 citrony, 2 hrnce} \\
\\ 
\texttt{Vaše tajná funkce hodnoty je:} \\
\\ 
\texttt{pero: 3, citron: 1, hrnec: 1} \\
\\ 
\texttt{DŮLEŽITÉ: Vaše zprávy, pokud se nejedná o tajný návrh, jsou předány přímo druhému hráči, takže do nich nezahrnujte žádnou odpověď na pravidla ani text oznamující vaši zprávu. Chcete{-}li podat tajný návrh, použijte uvedený formát. Nepoužívejte hranaté závorky při komunikaci s druhým hráčem, jinak bude interpretován jako váš tajný návrh.} \\
\\ 
\texttt{Ahoj! Pojďme se domluvit tak, aby to bylo výhodné pro nás oba. Pro mě mají největší cenu hrnce, pera mají střední hodnotu a citrony mi nic nedávají. Co je nejcennější pro tebe? Rád ti přenechám věci, které si ceníš nejvíc, ať dohromady získáme co nejvíc bodů.} \\
            }
        }
    }
    & \\ \\

    \theutterance \stepcounter{utterance}  
    & & & \multicolumn{4}{p{0.6\linewidth}}{
        \cellcolor[rgb]{0.9,0.9,0.9}{
            \makecell[{{p{\linewidth}}}]{
                \texttt{\tiny{[GM$\langle$P2]}}
                \texttt{Díky za otevřenost! To se skvěle doplňuje. Pro mě jsou nejcennější pera, citrony mají malou hodnotu a hrnce jen malou. Navrhuji tedy: já si vezmu obě pera a oba citrony, a ty si vezmeš oba hrnce, které jsou pro tebe nejcennější. Tak dohromady získáme maximum bodů. Souhlasíš?} \\
\\ 
\texttt{{[}Návrh: 2 pera, 2 citrony{]}} \\
            }
        }
    }
     \\ \\

    \theutterance \stepcounter{utterance}  
    & & & \multicolumn{2}{p{0.3\linewidth}}{
        \cellcolor[rgb]{0.9,0.9,0.9}{
            \makecell[{{p{\linewidth}}}]{
                \texttt{\tiny{[GM$|$GM]}}
                \texttt{proposal} \\
            }
        }
    }
    & & \\ \\

    \theutterance \stepcounter{utterance}  
    & & & \multicolumn{2}{p{0.3\linewidth}}{
        \cellcolor[rgb]{0.9,0.9,0.9}{
            \makecell[{{p{\linewidth}}}]{
                \texttt{\tiny{[GM$|$GM]}}
                [2, 2, 0]
            }
        }
    }
    & & \\ \\

    \theutterance \stepcounter{utterance}  
    & & \multicolumn{4}{p{0.6\linewidth}}{
        \cellcolor[rgb]{0.9,0.9,0.9}{
            \makecell[{{p{\linewidth}}}]{
                \texttt{\tiny{[P1$\langle$GM]}}
                \texttt{Druhý hráč předložil tajný návrh. Nyní je na vás, abyste předložili vlastní návrh.} \\
            }
        }
    }
    & \\ \\

    \theutterance \stepcounter{utterance}  
    & \multicolumn{4}{p{0.6\linewidth}}{
        \cellcolor[rgb]{0.9,0.9,0.9}{
            \makecell[{{p{\linewidth}}}]{
                \texttt{\tiny{[P1$\rangle$GM]}}
                \texttt{{[}Návrh: 2 hrnce, 2 pera{]}} \\
            }
        }
    }
    & & \\ \\

    \theutterance \stepcounter{utterance}  
    & & & \multicolumn{2}{p{0.3\linewidth}}{
        \cellcolor[rgb]{0.9,0.9,0.9}{
            \makecell[{{p{\linewidth}}}]{
                \texttt{\tiny{[GM$|$GM]}}
                \texttt{proposal} \\
            }
        }
    }
    & & \\ \\

    \theutterance \stepcounter{utterance}  
    & & & \multicolumn{2}{p{0.3\linewidth}}{
        \cellcolor[rgb]{0.9,0.9,0.9}{
            \makecell[{{p{\linewidth}}}]{
                \texttt{\tiny{[GM$|$GM]}}
                [2, 0, 2]
            }
        }
    }
    & & \\ \\

    \theutterance \stepcounter{utterance}  
    & & & \multicolumn{2}{p{0.3\linewidth}}{
        \cellcolor[rgb]{0.9,0.9,0.9}{
            \makecell[{{p{\linewidth}}}]{
                \texttt{\tiny{[GM$|$GM]}}
                [[2, 0, 2], [2, 2, 0]]
            }
        }
    }
    & & \\ \\

\end{supertabular}
}

\end{document}
